\section{A broad side of constant width}
\label{sec:a-constant-broad-high-neighborhood-case}

To turn Lemma~\ref{lem:bias-purity} into initial discrepancy, we need to
exclude available nearly pure patches with two small power widths. We
approach this through two broad-side constructions. The present
constant-width lemma handles an opposite-color alternative in the proof
of Lemma~\ref{lem:small-polynomial-broad-side}. The latter allows the
broad-side law to have a small polynomial density cap, or equivalently
logarithmic width. After stabilization, it supplies the lower bound on
both color densities in~\eqref{eq:source-1}, which the grid construction
of Section~\ref{sec:a-small-power-key-grid} uses to exclude these available
nearly pure patches.

\begin{lemma}[A broad side of constant width]
\label{lem:constant-broad-high-neighborhood}
Work in the counterexample sequence fixed above. Suppose for one color \(G\) there is a finite random tag \(I\) with laws \(\mu_I\) on \(X\) of width \(\le n^\gamma\) (\(0<\gamma<1\) fixed), \(\mathbb E\mu_I\le K'_0/N\) pointwise. For each \(I\), at least \(\chi N\) columns \(y\in Y\) have \(d_G(\mu_I,y)\ge .95\), with \(K'_0,\chi>0\) fixed. This yields a \(G\)-cube for large \(n\).
\end{lemma}

\begin{proof}
We construct probability rows on \(Y\) for the odd roles, assign them
injectively, and then construct rows on their common neighborhoods in
\(X\). The individual laws \(\mu_I\) may be too concentrated to serve a
constant fraction of the cube. Their balanced average will instead control
the mean loads of short random tuples in \(X\).

Each such tuple will be sampled to see several random columns. These
column requirements must vary across odd roles to spread their loads,
while only a small number of distinct requirements may occur around one
even role. We arrange this using counts and majority signs of disjoint
chunks of cube coordinates. Posterior distributions then relate three
objects: columns given streams of first-side labels, stream blocks given
those columns, and odd labels given the resulting blocks. All these
posteriors refer to specified raw experiments. We impose their numerical
bounds later, one sampling stage at a time.

\subsection{Parents, column indices and even types}
\label{sec:parents-macro-keys-and-even-types}

Call \(y\) good for \(I\) if \(d_G(\mu_I,y)\ge .95\), and define
\[
 (y,y')\in\mathcal H
 \quad\Longleftrightarrow\quad
 \Pr_I(y\text{ and }y'\text{ are good for }I)\ge\chi^2/2.
\]
The average of this probability over uniform \(y,y'\in Y\) is at least
\(\chi^2\), by squaring the size of each tag's good-column set. Hence
\(\mathcal H\) contains a positive constant fraction of all pairs.
Consequently a positive constant fraction of labels have a positive
constant fraction of partners in \(\mathcal H\). Draw \(V_0\) uniformly
from these labels. For each bin \(w\), to be defined below, independently
draw \(A_w\) uniformly from the partners of \(V_0\), conditional on
\(V_0\). Both the law of \(V_0\) and each conditional law of \(A_w\)
have atoms at most \(O(1)/N\). All draws use the retained labels under
consideration.

Fix constants
\[
 a_0=1/3<a_1<\cdots<a_8<\tau_0<\tau_1<\log2,
\]
and take \(\delta>0\) sufficiently small relative to their gaps. Given
the parents, draw a stream \(W_w\) at each bin, independently across bins
and in independent segments of a fixed large length \(q_0\). To draw one
segment, first choose a fresh tag conditioned on both parents being good,
and then draw \(q_0\) independent labels from \(\mu_I\) conditioned on
adjacency to both parents. The labels in a segment are independent
conditional on its tag. Each has simultaneous-hit probability at least
\(.90\) before this last conditioning. Thus the segment law has density
at most
\[
 (2/\chi^2)(.90)^{-q_0}\le e^{a_0q_0}
\]
relative to \(R'_0=\mathbb E_I\mu_I^{\otimes q_0}\), upon increasing the
fixed \(q_0\). For a length \(u\) divisible by \(q_0\), write
\(R[u]\) for the product of \(u/q_0\) such reference segments. Each of
its coordinate marginals is at most \(K'_0/N\), and its joint density
relative to product uniform measure is at most \(e^{u n^\gamma}\).
Only finite streams are needed. All prefix and block lengths below are
rounded upward to multiples of \(q_0\).

\paragraph{Counts, signs, and column indices.}
Partition some cube coordinates into \(d_c=300\) coarse chunks, each of
length \(\lfloor n^{1/5}\rfloor\). In each chunk, partition the possible
counts of ones into consecutive intervals, or bins, of probability at
most \(2n^{-.04}\), with \(O(n^{.04})\) endpoints. The binomial atom
bound \(O(n^{-.1})\) permits this choice. Let \(P\) be the vector of
counts and \(w(P)\) its vector of bins. The coarse key \(i(P)\) records
\(w(P)\) and a flag: it is interior if no single bit flip changes a bin,
and boundary otherwise. An interior key uses parent \(A_w\) and the
list \(C_i=\{w\}\); a boundary key uses parent \(V_0\) and the list
consisting of \(w\) and its grid neighbors. In particular,
\(w(P)\in C_{i(P')}\) whenever \(P'=P\) or \(P'\) results from one
coarse bit flip.

Use \(m=\lceil n^\alpha\rceil\) further disjoint chunks, each of the
same odd length asymptotic to \(n^{.3}\), where \(\alpha>0\) is a small
fixed constant. Their majority signs give \(t\in Q_m\). Let \(F\) be
the set of chunks whose count is at distance \(1/2\) from mid-weight;
these are exactly the signs that one bit flip can change. Let \(j\)
count the chunks whose count is at distance at most \(R_f=5.5\) from
mid-weight. Thus \(j\) counts chunks, not coordinates within them, and
its cutoff is independent of the coarse boundary flag. A role is low if
\(j\le J=\lfloor m^{1/20}\rfloor\), and high otherwise. All remaining
cube coordinates are residual. On either actual cube parity the signs
are uniform. Binomial atom bounds and a union over subsets of chunks give
\[
 \Pr(j\ge h)\le n^{-.13h}\quad(1\le h\le m),\qquad
 \Pr(i(P)\text{ is boundary})\le n^{-.05}.
\]

The hidden columns are indexed by
\(\ell=(i,t,j)\), \(0\le j\le J\), and by \(\ell=(i,*)\) in the
high case. An odd role uses the index determined by its counts and signs.
Given the base data \((V_0,A,W)\), independently draw \(U_\ell\) for
all these indices. A low \(U_\ell\) is one column; a high \(U_\ell\)
is a tuple of \(s=\lceil K_sJ\log m\rceil\) independent columns.
Write \(s_\ell=1\) or \(s\) for its length. Its coordinate law
\(\pi_\ell\) is the following posterior in the raw parent-and-stream
experiment:
\begin{itemize}
\item at an interior key, the law of \(A_w\) given
\(V_0,W_w[1:b_j]\);
\item at a boundary key, the law of \(V_0\) given
\((A_{w'},W_{w'}[1:b_j])_{w'\in C_i}\).
\end{itemize}
Here
\[
 u_j=\lceil K_1(j+4)\log m\rceil,\qquad b_j=u_{j+1},
\]
with segment rounding, and high keys use \(j=J\). We draw these variables
also at indices that occur only in the padded lists defined next.

\paragraph{Even types and the lengths of their arrays.}
For a low even role, with coarse data \(P,w\), let \(S\) consist of the
distinct keys
\begin{itemize}
\item \((i(P'),t,j)\), where \(P'=P\) or is obtainable by one coarse
flip;
\item \((i(P),t^h,j)\), \(h\in F\), and
\((i(P),t,j\pm1)\) when in range, replacing severity \(J+1\) by the
high key. Here \(t^h\) is \(t\) with sign \(h\) reversed.
\end{itemize}
Its type is \(K=(w,S,j)\). At a high even role take
\(K=(w,S,*)\), where \(S\) contains the high keys at the same coarse
range of \(i(P')\)'s. If \(j=J+1\), there is also an optional low
neighbor key \((i(P),t,J)\). This optional key is excluded from both
\(S\) and the type \(K\). These lists cover the keys at all actual
neighbors, since a sign flip leaves \(j\) unchanged. Every key paired
with \(w\) has \(w\in C_i\).

We now fix the lengths and indices of the arrays; their block law
\(P_K\) will be defined in Step~2. At each possible center ID, arrays of
different types will be independent conditional on the base and hidden
columns. For a low type, blocks have length \(u=u_j\), and the array
has total length
\[
 k_j=u_j\left\lceil K_2\big((j+4)\log m+m^{.02}\big)/u_j\right\rceil.
\]
For a high type, put \(u=u_*=\lceil\eta\log m\rceil\), with segment
rounding, and \(k_*=u_*\lceil m^{.005}/u_*\rceil\). Its full pool has
\(\lceil e^{K_hu_*}k_*/u_*\rceil\) blocks, of total length \(M_*\).
If an optional low key is designated, the used tuple is formed from the
first \(k_*/u_*\) blocks all of whose entries hit that column. Without
an optional key it consists of the first \(k_*/u_*\) blocks. At fixed base, the full
pool law will use only the keys in the type \(K\), so it is independent
of the optional column's value before this extraction.

A tuple reference \(c\) specifies its ID and type, and in the high case
also the block index subset. Its length \(k_c\) is \(k_j\) or \(k_*\).
The constants \(K_h,K_1,K_2,K_s\) are fixed and large; after \(K_h\)
choose \(\eta\) so small that \(M_*=o(m^{.03})\). We give the remaining
order of choices below, taking \(\alpha\) sufficiently small last.
Undefined numerical formulas on unused or failure branches use fixed
fallback values depending only on their declared local inputs.

The raw sampling order is
\[
 V_0\ \longrightarrow\ (A_w,W_w)_w\ \longrightarrow\ (U_\ell)_\ell
 \ \longrightarrow\ \text{center arrays}.
\]
Prospective positions, activations and choice ties are sampled separately;
an array is generated even at an ID whose center is absent. Each posterior
uses this raw experiment and only the observations explicitly listed in
its definition. When a candidate value changes, every later raw kernel
depending on it is evaluated at that new value. In particular, unobserved
arrays are integrated under their new laws. Successful-history conditions
are included only when explicitly retained as local gates.

\subsection{Predictive estimates at the first two levels}
\label{sec:predictive-estimates-at-the-first-two-levels}

\paragraph{Step 1: deleting a stream prefix from a parent posterior.}
For a key \(\ell\), remove the observation \(W_w[1:u]\) from the
specified definition of \(\pi_\ell\), retaining its other observations,
and call the resulting posterior \(\pi_{\ell,-w,u}\). The shift
\(b_j=u_{j+1}\) ensures that each prefix needed for a listed or optional
key lies within the observed stream. Given the retained observations,
the incoming prefix has likelihood at most \(e^{a_0u}\) relative to
\(R[u]\), for every parent value in support: it comprises complete
segments, and the partner is among the retained data.

The marginal density of this prefix is smaller than \(e^{-\delta u}\)
on a set of raw probability at most \(e^{-\delta u}\). Off that set,
Bayes' formula and the gap \(a_1-a_0\) give the test
\[
 \pi_\ell(y)\le e^{a_1u}\pi_{\ell,-w,u}(y)
 \qquad\text{for every }y.
\]
We also require
\(N\max\pi_\ell\le\exp(K'b_j)\), for a sufficiently large fixed
\(K'\). Its raw exception is at most \(e^{-\delta b_j}\). At an
interior key this follows from the full stream likelihood conditional
on \(V_0\). At a boundary key use product uniform measure for the
bounded list \((A_{w'})_{w'\in C_i}\) and the segment references for
their streams. The likelihood for each candidate \(V_0\) is at most
\((O(1)e^{a_0b_j})^{|C_i|}\). The prior atom bound and the same
small-denominator estimate give the asserted cap.

\paragraph{Step 2: reconstructing a stream block from hidden columns.}
Fix a type \(K\) and all base data except \(W_w[1:u]\). Its conditional
raw law \(P_0\) is the product of the true-parent segment laws. Let
\(\mathcal G_K\) be the set of alternative blocks \(z\) for which the
Step~1 comparisons at \((w,u)\) hold for every key in \(S\). At a high
type, also require them for all potentially used optional keys. This
involves all coarse parts represented in \(S\), but not the signs:
the low \(j=J\) and high priors have the same coordinate law. In
particular, \(\mathcal G_K\) is defined without observing any \(U\).

Set \(\rho_\ell=\pi_{\ell,-w,u}^{\otimes s_\ell}\). Relative to
\(\rho_\ell\), let \(L_\ell(z;\theta_\ell)\) be the likelihood of
\(U_\ell=\theta_\ell\) when the base block is replaced by \(z\).
Off the support of \(\rho_\ell\) the gated likelihood can be set to
zero. Relative to the product reference \(\prod_{\ell\in S}\rho_\ell\),
the observations \(U_S\), with the event
\(W_w[1:u]\in\mathcal G_K\) retained, have subdensity
\[
 m_S(\theta)=\int1_{\mathcal G_K}(z)
       \prod_{\ell\in S}L_\ell(z;\theta_\ell)\,dP_0(z).
\]
Define \(P_K\) by normalizing this integrand in \(z\) at the observed
\(U_S\). Dropping the factor for \(\ell\), but retaining the same
gate, defines \(m_{S-\ell}\) and \(P_{K,-\ell}\).

We require positivity and the denominator bounds
\[
 m_S\ge e^{-\delta u\sum_Ss_\ell},\qquad
 m_S/m_{S-\ell}\ge e^{-\delta us_\ell}.
\]
On the true-block gate, each exception has raw probability at most its
threshold. For the ratio test, for example,
\[
 \int_{\{m_S<\varepsilon m_{S-\ell}\}}m_S
       \,d\Bigl(\prod_{\ell'\in S}\rho_{\ell'}\Bigr)
 \le\varepsilon\int m_{S-\ell}
       \,d\Bigl(\prod_{\ell'\in S}\rho_{\ell'}\Bigr)
 \le\varepsilon.
\]
The last integral is at most one, and all unlisted keys have simply been
integrated out. On the true-block gate \(m_S>0\) almost surely, which
also implies \(m_{S-\ell}>0\). For a high type, adding one optional low
key gives the analogous test
\(m_{S+\ell}/m_S\ge e^{-\delta u_*}\), with the same probability
bound.

Combining these lower bounds with Step~1 gives
\[
 \frac{dP_K}{dP_{K,-\ell}}\le e^{a_2us_\ell},\qquad
 \frac{dP_K}{dR[u]}\le A_K^u,\qquad
 A_K=\exp\bigl(K''(1+\sum_{\ell\in S}s_\ell)\bigr),
\]
where \(K''\) is a sufficiently large fixed constant. The first
comparison holds also at any replacement value of \(\theta_\ell\)
passing its ratio test. We can now generate each center array as
independent \(P_K\)-blocks of the lengths already specified.

The support of this posterior enforces the required adjacencies. Every
\(z\) in the support of \(P_0\) leaves the base data with positive raw
probability. If a listed column has positive posterior likelihood at
that base, the stream-generation rule says that it sees every entry of
\(z\). For an optional column the added likelihood is at most
\(e^{a_1u_*}\), whereas its expectation under \(P_K\) is
\(m_{S+\ell}/m_S\ge e^{-\delta u_*}\). Its positive-likelihood set
therefore has \(P_K\)-mass at least \(e^{-(a_1+\delta)u_*}\), and every
block in that set hits the optional column. Choosing \(K_h\) large
makes the probability of too few such blocks in a high pool \(o(1)\),
uniformly at histories passing these tests.

\subsection{State neighborhoods and odd predictive estimates}
\label{sec:state-neighborhoods-and-odd-predictive-estimates}

We group actual cube roles into states that retain their residual bits,
exact coarse counts, and the count in every majority-sign chunk, with one
exception: counts at distances \(5.5\) and \(6.5\) on the same side of
mid-weight are merged. A state also records the exact value of \(j\).
Encode each count field and the \(j\)-field by a one-hot vector. Together
with the residual bits this embeds the states in
\(Q_d\), \(d=(1+o(1))n\).

A state determines parity despite the merged counts. Give each merged
symbol its outer representative, at distance \(6.5\), as a baseline.
Replacing it by its inner representative changes the count by one and
reverses parity. The number of replacements is \(j\) minus the number
of unmerged counts within the cutoff. This determines the fine-coordinate
parity; the coarse counts and residual bits are exact. It also determines
the key, type and optional designation. The exact coarse counts here are
more information than will be retained in
Lemma~\ref{lem:small-polynomial-broad-side}.

Join an odd state to an even state when some actual cube edge realizes
that pair. Each state has degree \(O(n)\), and any two even states
adjacent to one odd state have ambient distance at most \(D=8\). Indeed,
one bit flip changes only boundedly many one-hot entries, including the
possible change of \(j\). A flip changing \(j\) passes between distances
\(5.5\) and \(6.5\). Thus all low neighbors of a high odd state, when
present, form one state; the same holds for the high neighbors of a low
state. The padded lists above cover all this incidence.

Use even states as height sites and choose one center ID at each, shared
by all actual roles represented there. An ID is a location-level pair in
the ambient cube. We aim to use at most \(T=\lceil m^{.001}\rceil\)
distinct IDs around each odd state. For the height device take
\(r=\lfloor\rho n\rfloor\), \(\lambda=n^{10}\), and parameters with
\(2n^b\le T\) and \(\theta>.9\). Its exponents are chosen after
\(\alpha\), with \(b>b_0>0\) much smaller than \(.001\alpha\), and
\(\sigma,\zeta,1-\theta\) small enough for its inequalities. These
height exponents are unrelated to the stream lengths \(b_j\). Choose
\(\rho>0\) so small that the binary ball entropy at \(2\rho\), in
natural logs, is less than \(\log2-\tau_1\). We will also use the short
height computation of radius \(D\lfloor\sqrt m\rfloor\).

\paragraph{Observation records and their counts.}
Fix an odd state \(b\) and a mapping of its neighboring even states to
at most \(T\) IDs. Its observation record lists distinct ID/type arrays,
so repeated uses of the same array are recorded only once. At a low state
we observe every used full low tuple and, if a high neighbor is present,
its full pool and the block subset selected to hit \(U_\ell\). Call
that subset the mask. At a high state we observe all used full high pools
and the possible one used low tuple. We also record which high
type/optional-key combinations are needed at each ID, including the
absence of an optional key. Their block subsets are computed from the
observed pools and optional columns. If there are too few hits, the test
uses the first blocks as a fallback; actual usable tuples will be ensured
separately. A tuple reference \(c\) is therefore a specified set of
primitive blocks in one of the observed arrays.

There are \(O(T+j)\) low tuples around a low state. To see this, low types
with unchanged \(t,F\) have boundedly many generic variants, determined
by coarse incidences and severity; each variant uses a subset of the ID
list. Every other variant changes a sign or \(F\), and hence comes from
a chunk at distance at most \(1.5\) from mid-weight. There are at most
\(j\) such chunks and constantly many relevant flips per chunk. Each
of their \(O(j)\) neighboring states uses one ID. Ambiguous merged
chunks do not change signs or \(F\), so they produce only generic
variants regardless of how many there are.

At a high state, there are \(O(T)\) full pools of boundedly many types.
Optional choices vary only near the interface, where \(j\le J+2\).
Their generic variants retain the sign, while the other variants arise
from \(O(J)\) relevant neighboring states. Thus there are only
\(O(T+J)\) needed high tuple references, even though several may use
subsets of the same pool.

The counts of actual records below assume that the \(r\)-ball of each
even state adjacent to \(b\) has at most \(2\lambda\) prospective centers
at every level, as will be required in the local validity event.
At fixed center presence, this leaves polynomially
many possible IDs. Choosing their lists and assigning the exceptional
states gives logarithmic record counts at most
\[
 \begin{array}{ll}
 O(T\log n+(j+1)\log T+k_*1_{j=J}),&b\text{ low},\\
 O(T\log n+(J+1)\log T),&b\text{ high}.
 \end{array}
\]
\Needspace{3\baselineskip}
The low mask contributes at most
\(\log\binom{M_*/u_*}{k_*/u_*}=O(k_*)\). At high states we count
optional-key assignments, then compute their subsets; we do not enumerate
all subset combinations. If only the full observation list is recorded,
its logarithmic count is \(O(T\log n)\).

For history tests, before positions have been sampled, use abstract
records with the IDs renamed in \([T]\). Group them by coarse bin,
central sign and severity. Their logarithmic counts are
\[
 \begin{array}{ll}
 O(T\log T+(j+1)\log m+k_*1_{j=J}),&b\text{ low},\\
 O(T\log T+(J+1)\log m),&b\text{ high}.
 \end{array}
\]
Besides generic variants, a low record specifies \(O(j+1)\) chunk
indices and their finitely many near-center statuses. A high record
needs such information only at the interface. Coarse data enter through
a bounded neighborhood of bins. In particular, regular Step~2 low-type
patterns number \(m^{O(j+1)}\); the small-prefix \(u_*\) tests have
boundedly many patterns per bin modulo the sign of their one optional
key. These counts distinguish the full arrays observed, the tuple
references extracted from them, and the records used to enumerate tests.

\paragraph{Step 3: reconstructing an odd label from center arrays.}
Fix the base and every hidden column except the target
\(\vartheta=U_\ell\) at \(b\). Its raw prior is
\(\pi_\ell^{\otimes s_\ell}\). At a low state, condition also on the
possible high pool, whose law does not use this optional low key. The
remaining observations are the low-array blocks. At a high state, the
observations instead comprise every full high pool and the possible one
low tuple. In either case these observed blocks are independent given
\(\vartheta\), with laws \(P_K(\cdot\mid\vartheta)\); their types
have \(\ell\in S\).

For each observed block use \(P_{K,-\ell}\) as a reference law. Let
\(g_e(\vartheta)\) be the likelihood of block \(e\) relative to this
law. The candidate gate \(\Theta(\vartheta)\) requires positive
\(m_{S-\ell}\) and the Step~2 lower ratio \(m_S/m_{S-\ell}\) for every
observed type. These conditions use no array entries. At a low state
with a high pool, also include the indicator that \(\vartheta\) selects
the specified successful mask from that fixed pool. There is no mask
indicator at a high state. Thus, with \(c\) a same-mode tuple reference,
\[
 M=\int\Theta(\vartheta)\prod_e g_e(\vartheta)
       \,d\pi_\ell^{\otimes s_\ell}(\vartheta),\qquad
 M_{-c}=\int\Theta(\vartheta)\prod_{e\notin c}g_e(\vartheta)
       \,d\pi_\ell^{\otimes s_\ell}(\vartheta).
\]
At low states the products use only low blocks; at high states they use
all the observed blocks, and deletion may remove a strict subset of one
pool. On the candidate gate each block density is at most
\(e^{a_2|e|s_\ell}\), where \(|e|\) is its length.

Let \(k'_j\) be the minimum \(k_l\) with \(|l-j|\le1\). Require
positivity and
\[
 M/M_{-c}\ge e^{-\delta k_cs_\ell},\qquad
 M\ge
 \begin{cases}
 e^{-\delta k'_j},&b\text{ low},\\
 e^{-\delta s},&b\text{ high}.
 \end{cases}
\]
For a fixed observation record and fixed subsets, each exception
intersected with the
true-target gate costs at most its threshold in the raw joint experiment.
This is the subdensity calculation of Step~2, now using
\(\int M_{-c}\le1\). At low states it includes occurrence of the
specified successful mask. At high states an individual ratio test can
be unioned over all subsets of one pool at cost \(\exp(O(k_*))\),
which is smaller than its \(\delta k_*s\) exponent. No union over all
joint choices of subsets is needed.

Let \(\mathcal P\) and \(\mathcal Q_c\) be the normalized full and
deleted target integrands. Their gates do not include the denominator
tests just imposed. At a low state, \(\mathcal P\) is supported on
labels adjacent to all needed tuples, including the high tuple specified
by the mask. The block likelihood and denominator bounds give
\(\mathcal P\le e^{a_3k_c}\mathcal Q_c\). We reserve a further
multiplier \(e^{\delta k'_j}\) for selection, still below
\(e^{a_4k_c}\). Even after that multiplier the normalized atom cap is
\(e^{D_L}\), where \(D_L=m^{.15}\): its logarithm costs only
\(O(b_j+(T+j)\max_{|l-j|\le1}k_l)\). The independent high pool is not
part of this likelihood cost.

At a high state the bounds are
\[
 \log(\mathcal P/\mathcal Q_c)\le a_3k_cs,\qquad
 \log(N^s\max\mathcal P)\le K_0sJ\log m.
\]
There are \(O(TM_*)=o(J)\) full high-pool entries and at most one low
tuple of length \(O(J\log m)\); together with the prior and denominator
bounds these give the second inequality. The fixed \(K_0\) may depend
on \(K_1,K_2\), but not on \(K_s\).

\paragraph{High rows: a common law for all deletion costs.}
For a high row we need one law on index--label pairs that has small
expected likelihood cost against every needed high deletion reference.
Fix the mapping and observed data in the true-target-gated raw model,
with the preceding denominator tests satisfied. The remaining target \(\vartheta\) has law \(\mathcal P\). The needed
block subsets are now fixed by the full pools and the other, optional
low keys; they do not depend on the remaining target draw. Let
\(r_{\rm ref}=O(T+J)\) be their number. If it is zero, the cost
constraints below are vacuous: restrict the index--label law to the
density-good set defined below and normalize, without using a grid.

For a supported true prefix \(\vartheta_{<h}\), let
\(\mathcal P_h\) and \(\mathcal Q_{c,h}\) be the next-coordinate
conditionals. They are functions of this prefix. Joint domination ensures
that the deletion conditional is defined there. Smooth it with a small
fixed uniform component to obtain \(\widetilde{\mathcal Q}_{c,h}\),
and put
\[
 E_h^c(y)=k_c^{-1}
   [\log(\mathcal P_h(y)/\widetilde{\mathcal Q}_{c,h}(y))]_+.
\]
If \(Y\) has law \(p\), then
\(\Pr_p(\log(q(Y)/p(Y))>z)\le e^{-z}\): on that set
\(p(Y)<e^{-z}q(Y)\). The positive part of this negative log has
exponential moment at parameter \(1/2\) at most \(2\). Applying this
conditionally at successive true prefixes shows that the sum of negative
parts is \(O(s)\), except with probability \(e^{-\Omega(s)}\) per
reference. By the chain rule, the sum of the full unsmoothed logs is at
most \(a_3k_cs\). Smoothing increases a positive log by at most a
fixed constant. Consequently
\[
 \frac1s\sum_{h=1}^s E_h^c(\vartheta_h)\le a_3+o(1).
\]
The same argument against uniform measure gives
\[
 \frac1s\sum_{h=1}^s
 \frac{[\log(N\mathcal P_h(\vartheta_h))]_+}{D_H}
 \le K_0/K_D+o(1),\qquad D_H=K_DJ\log m.
\]
We choose \(K_D\) large enough to make this last bound a small fixed
constant.

These are averages at the labels actually drawn. To control the averages
of the conditional laws, fix a convex weight vector \(\xi\) on the
references and a large fixed \(L_0\). For the filtration
\(\mathcal F_h=\sigma(\vartheta_1,\ldots,\vartheta_h)\), conditional
on the fixed observations, set
\[
 C_h^\xi(y)=\min\{L_0,\sum_c\xi_cE_h^c(y)\},\qquad
 D_h^\xi=C_h^\xi(\vartheta_h)
                -\sum_y\mathcal P_h(y)C_h^\xi(y).
\]
The \(D_h^\xi\) are martingale differences bounded in absolute value
by \(L_0\). On the preceding actual-path bounds, failure of
\[
 \frac1s\sum_h\sum_y\mathcal P_h(y)C_h^\xi(y)
       \le a_3+\delta'
\]
forces \(\sum_hD_h^\xi<-\delta's/2\), for large \(n\). Its probability
is at most \(\exp(-\Omega_{\delta',L_0}(s))\). This concentration is
applied before conditioning on any successful path. Applying it also to
\(\min\{1,[\log(N\mathcal P_h(y))]_+/D_H\}\) bounds the average
conditional mass of labels with \(N\mathcal P_h(y)>e^{D_H}\) by a
small fixed constant \(\eta_D\).

Here is a finite set of directions sufficient for all the cost
constraints. Fix a small \(\varepsilon_g>0\), and for
\(\xi\) in the reference simplex set
\[
 a_c=\lceil r_{\rm ref}\xi_c/\varepsilon_g\rceil,\qquad
 A=\sum_c a_c,\qquad \widehat\xi_c=a_c/A.
\]
Then \(A\le(1+\varepsilon_g^{-1})r_{\rm ref}\) and
\(\xi_c\le(1+\varepsilon_g)\widehat\xi_c\) for every \(c\).
The number of nonnegative integer vectors with this bound on their sum
is \(\exp(O_{\varepsilon_g}(r_{\rm ref}))\), by the composition count
\(\binom{\lfloor(1+\varepsilon_g^{-1})r_{\rm ref}\rfloor+r_{\rm ref}}
{r_{\rm ref}}\). Thus all the predictable bounds hold on this grid,
with total failure \(e^{-c's}\) for a fixed \(c'>0\), since
\(T+J=o(s)\). This is a conditional bound at fixed data; integrating
against the gated data subdensity and adding the earlier denominator
exceptions gives the same form of bound for each mapping. The later
union over abstract mappings has a different count and will determine
how large \(K_s\) must be.

Now fix a path passing these tests, and define the probability law
\(\nu(h,y)=\mathcal P_h(y)/s\) and set
\(G_D=\{(h,y):N\mathcal P_h(y)\le e^{D_H}\}\). For each grid direction
restrict \(\nu\) to \(G_D\) and to
\(\sum_c\widehat\xi_cE_h^c(y)\le L_0\). The removed mass is at most
\[
 \beta=\eta_D+(a_3+\delta')/L_0.
\]
After normalization, the unclipped expected cost in that direction is
at most \((a_3+\delta')/(1-\beta)\). Choose the small constants and
\(L_0\) so that \((1-\beta)^{-1}\le1+\varepsilon'\), with
\(\varepsilon'>0\) small. All these restricted laws belong to the same
compact convex set
\[
 \mathfrak F=\left\{R\ge0:\ \sum_{h,y}R(h,y)=1,\quad
 R(h,y)\le(1+\varepsilon')\nu(h,y)1_{G_D}(h,y)\right\}.
\]
Nonnegative costs and the coordinatewise grid bound show that for every
convex direction \(\xi\), some member of \(\mathfrak F\) has weighted
expected cost at most \(a_3\) plus a chosen small error. If no member
satisfied all the individual bounds, separation of the compact convex
set of cost vectors from the region below these bounds would give a
nonnegative separating vector. After normalizing it, this would contradict
the directional conclusion. Hence one law \(\mathcal R\in\mathfrak F\)
satisfies all the cost bounds.

Its support hits every required tuple: a positive \(\mathcal P_h(y)\)
has a positive posterior continuation, all of whose coordinates hit the
observed blocks. Moreover,
\[
 \left[\log\frac{\mathcal R(h,y)}
   {\widetilde{\mathcal Q}_{c,h}(y)/s}\right]_+
 \le \log(1+\varepsilon')+k_c E_h^c(y).
\]
The gap from \(a_3\) to \(a_4\) therefore yields
\[
 \mathcal R(h,y)\le 2e^{D_H}/(sN),\qquad
 \mathbb E_{\mathcal R}\!\left(
  [\log(\mathcal R(h,y)/(\widetilde{\mathcal Q}_{c,h}(y)/s))]_+
                         \right)\le a_4k_c.
\]
Choose one such law deterministically from the local data. A deleted
reference uses only the full observation list with those primitive blocks
removed, together with the fixed history and true prefixes. It does not
use other optional subset assignments. True prefixes are legitimate fixed
data here because the hidden columns were sampled before the arrays.
The clip parameters can be chosen with a rate \(c'>0\) independent of
\(K_s\). The path tests include positivity and true support; these hold
almost surely in the true gated model.

\subsection{Making history tests simultaneous}
\label{sec:making-history-tests-simultaneous}

We next sample the base and hidden columns so that Steps~1 and~2 hold
throughout. At the resulting history, every abstract mapping must also
have Step~3 failure probability over independent center arrays at most
\[
 e^{-c_Lk'_j}\quad\text{at low states},\qquad
 e^{-c_Hs}\quad\text{at high states},
\]
for fixed small \(c_L,c_H>0\). We achieve this in five stages, restricting
only the newly drawn variables at each entering history. Later avoidance
probabilities never reweight earlier histories.

For Step~2 a failure always means the failed numerical test intersected
with its true-block gate \(W_w[1:u]\in\mathcal G_K\). For Step~3 it
means failure at the specified record intersected with its true-target
gate, at a base passing Step~1. In a low record with a specified high
mask, the event includes occurrence of that successful mask. The earlier
stages will enforce the block and key-ratio gates; usable hits for selected
centers are enforced below. These definitions keep raw subdensity bounds
as intersection bounds, rather than conditioning them on future success.

A regular Step~1 or~2 test has scale \(L=u_j\), including the prior cap
at severity \(j\), with the high prior cap assigned \(j=J\). A
small-prefix test has scale \(L=u_*\). Each raw failure probability is
at most \(e^{-\delta L}\). The regular pattern count per bin and
central sign is \(m^{O(j+1)}\), so a bound such as
\(e^{-\delta u_j/16}\) sums to \(o(1)\) over all regular patterns and
severities if \(K_1\) is large. Small-prefix tests have boundedly many
patterns modulo sign translation, and
\(e^{-\delta u_*/16}\to0\).

An alarm at an entering history \(H\) is an event
\[
 \left\{\Pr_{\rm raw\ future}(B\cap\text{local gate}\mid H)>q\right\}.
\]
If the expectation of this conditional probability is at most \(p\),
Markov bounds the alarm probability by \(p/q\). Its scope comprises the
current-stage variables occurring in the canonical record: base bins in
the relevant \(C_i\)-lists, hidden keys in the relevant type lists, and
any specified optional keys. Future array variables are integrated out.
This definition determines the scopes used in each local lemma below.

For these alarm families, the gated likelihoods are evaluated only where
their defining ratios have positive denominators. Auxiliary continuations
may still reach undefined branches. Fix their fallback rules invariant
under the bin, sign and ID renamings used in the abstract records: for
example, fix an order on retained host labels, use the first label and
repeated copies when a probability law is undefined, and fixed numerical
values for undefined scalar formulas. These choices depend on no
additional keys or arrays. Thus the symmetry assertions below hold also
for the completed experiments.

\paragraph{Conditioning stage 1: the global parent.}
For each abstract Step~1 or~2 pattern, the conditional failure probability
given \(V_0\) is a function of this one draw. Its mean is at most
\(e^{-\delta L}\). Restrict \(V_0\) so that every such function is at
most \(e^{-\delta L/2}\); the alarm probability per pattern is at most
\(e^{-\delta L/2}\).

The pattern unions just given apply globally here. Conditional on \(V_0\),
the bin variables are iid, every formula reads a bounded-radius bin
neighborhood, and only boundedly many coarse incidence shapes occur.
Likewise translating all signs does not change a raw failure probability:
the low keys being integrated have independent laws that do not depend
on their signs. No factor \(2^m\), or factor for all bin names, is
needed. The total excluded parent mass is \(o(1)\). The retained draw
does not change any of the posterior definitions.

\paragraph{Conditioning stage 2: the coarse base.}
Fix the selected \(V_0\) and start from the independent variables
\((A_w,W_w)\). Exclude Step~1 failures outright. Also exclude the
alarms that a Step~2 gated failure, averaged over its future hidden keys,
has probability greater than \(e^{-\delta L/4}\). Stage~1 bounds their
expectations by \(e^{-\delta L/2}\), so the Markov cost is at most
\(e^{-\delta L/4}\) per alarm.

Group the events at each bin. Their scopes are the base variables in the
bounded union of the record's \(C_i\)-lists. The dependency degree is
bounded, and the pattern sums are \(o(1)\). The conditional avoidance
lemma therefore applies with charges \(o(1)\). It gives a base passing
Step~1 and the stated conditional Step~2 bounds, separately at each
selected \(V_0\).

\paragraph{Conditioning stage 3: the high hidden keys.}
Given a successful base, the high \(U\)'s are initially independent.
The one-target calculations of Step~3 give, for any fixing of the other
keys, raw gated exceptions at most
\(e^{-c_{L0}k'_j}\) or \(e^{-c_{H0}s}\), with fixed small positive
constants \(c_{L0},c_{H0}\). These estimates integrate the target and
arrays. A high record specifies type/optional-key appearances and
computes their subsets from the pools; it does not condition on their
extraction in advance.

Restrict the high keys to pass high-only Step~2 tests. For tests involving
low keys, require their conditional raw failure probabilities to be at
most \(e^{-\delta L/8}\). Stage~2 and Markov give alarm probabilities
at most \(e^{-\delta L/8}\). Also require each high Step~3 failure,
still averaged over independent low keys and arrays, to have probability
at most \(e^{-c_{H0}s/2}\). Its alarm costs at most
\(e^{-c_{H0}s/2}\).

These events read only high keys in nearby coarse bins, hence have
bounded dependency degree after grouping by bin. The high Step~3
abstract count is \(\exp(O(T\log T+J\log m))\). Taking \(K_s\) large
makes its product with \(e^{-c_{H0}s/2}\) tend to zero. All grouped
probabilities are thus \(o(1)\), and another bounded-degree local lemma
provides the high-key history. The abstract alarms retain the sign and
ID symmetries; no later spatial selection is used in them.

\paragraph{Conditioning stage 4: separate optional pretrims.}
Given this base and high history, separately restrict each low \(U_\ell\)
to satisfy its optional small-prefix tests
\(m_{S+\ell}/m_S\ge e^{-\delta u_*}\). Only that low variable remains
random in such a test, and each is involved in boundedly many coarse
patterns. Their probabilities from Stage~3 are \(o(1)\), so the low
variables remain independent with a density multiplier \(1+o(1)\) per
variable.

A regular Step~2 pattern uses \(O(j+1)\) low keys, so its probability
multiplies by \((1+o(1))^{O(j+1)}\) and remains at most
\(e^{-\delta u_j/16}\). A high Step~3 pattern reads only \(O(J)\)
low keys: the possible low tuple's \(S\)-list and the optional keys
named in the record. Its bound remains \(e^{-c_{H0}s/3}\). For low
Step~3 the one-target bound holds at every fixing of the other low keys;
only the target's own distortion matters. Its bound remains
\(e^{-c_{L0}k'_j/2}\).

\paragraph{Conditioning stage 5: the low hidden keys.}
Start from the product of these trimmed low laws. Exclude all remaining
Step~2 failures. For Step~3 impose the conditional array-failure bounds
\(e^{-c_Lk'_j}\) and \(e^{-c_Hs}\), taking, for example,
\(c_L=c_{L0}/4\) and \(c_H=c_{H0}/6\). Markov applied to the Stage~4
bounds gives alarms of these same exponential orders.

Group by central sign, coarse bin and severity. The current variables
in a low group's scope are exactly the low keys in its listed types and
optional entries. They have bounded sign and severity distance and lie
in nearby bins. High groups involving low variables occur only near the
interface and read the optional keys and the possible low tuple's
\(S\)-list; high requirements with no low variables already hold.
Both the dependency degree and the number of groups touching one low
variable are polynomial in \(m\), with fixed exponents.

After the pattern unions, each group can have probability at most
\(m^{-K_{\rm deg}}\) for any prescribed fixed \(K_{\rm deg}\): choose
\(K_1\) large for the regular tests, then \(K_2\) large for the low
Step~3 logarithmic count
\(O(T\log T+(j+1)\log m+k_*1_{j=J})\), and \(K_s\) large for the
high count. Polynomially small charges with total \(o(1)\) at each
variable give the final local-lemma restriction. This completes the
construction of the key history and its uniform conditional bounds over
independent center arrays.

\paragraph{Order of constants.}
First fix \(q_0\), the gaps between the \(a_i\)'s, \(\delta\), and the
path clip parameters other than \(K_D\). Choose \(K_h\) for optional
block hits and then \(\eta\) so small that
\(K_h\eta+.005<.03\). The coefficients in the history pattern counts
are fixed before choosing \(K_1,K_2,K_s\). Next take \(K_1\) large for
regular history tests, then \(K_2\) large for low Step~3, then \(K_D\)
large relative to \(K_0\), and finally \(K_s\) large for the high
pattern unions. Increasing \(K_2\) or \(K_s\) need not decrease the
respective Step~3 exponential rates before counting patterns.

The scales give
\[
 TM_*=o(J),\qquad T\log n=o(k_*),\qquad
 T\log n=o(m^{.02}),\qquad T+J=o(s).
\]
Allow enough slack in the constants that each actual low configuration
count is smaller than its required fixed margin times \(k'_j\), and
that each high testing count is smaller than its required margin times
\(s\). These relations hold for every fixed \(\alpha>0\); choose
\(\alpha\) sufficiently small after all the likelihood constants,
including the following truncation constant.

\paragraph{Prior-heavy labels.}
For a block law \(P_K\), define its average coordinate marginal by
\(\bar P_K(x)=u^{-1}\sum_{i=1}^uP_K(z_i=x)\). Call a label
prior-heavy for \(K\) if \(N\bar P_K(x)>B_K\), where
\(B_K=A_K^{K_B}\). Fix a small \(\upsilon_0>0\) with
\(\upsilon_0+\tau_0/\tau_1<1\). At fixed successful key history and a
fixed tuple reference \(c\), the probability that more than
\(\upsilon_0k_c\) entries are prior-heavy is at most
\[
 2^{k_c}A_K^{k_c}
       \left(q_0K'_0/B_K\right)^{\upsilon_0k_c/q_0}.
\]
Indeed there are at most \(N/B_K\) heavy labels. A reference segment
hits this set with probability at most \(q_0K'_0/B_K\), and that many
heavy entries require at least \(\upsilon_0k_c/q_0\) segments to hit
it. The segments of \(R[k_c]\) are independent; the block domination
contributes \(A_K^{k_c}\), and the subset union contributes at most
\(2^{k_c}\). Choose \(K_B\) large enough that this is \(o(1)\), even
after the union over high subsets.

The following are instead deterministic averages over actual cube roles:
for sufficiently small \(\alpha\),
\[
 \mathbb E_{v\in A}B_{K(v)}\exp(O(k_*1_{j(v)>J}))=O(1),\qquad
 \Pr_{v\in B}(j(v)>J)e^{D_H}=o(1).
\]
At low severity \(\log B_K=O(j+1)\), except for an extra \(O(s)\)
at \(j=J\); at high severity it is \(O(s)\). The rarity bound
\(\Pr(j\ge h)\le n^{-.13h}\) outweighs these costs once \(\alpha\)
is small. The constants multiplying \(k_*\) are already fixed. Also
choose \(\alpha\) so that the Step~1 low-prior caps have logarithms
at most \(.02(j+1)\log n\).

\subsection{Height choices and the odd kernels}
\label{sec:height-choices-and-the-odd-kernels}

Fix a successful history through the hidden columns. Generate prospective
center presence and independent arrays according to their raw laws. At
an even site-level, first remove any center in its \(r\)-ball whose
tuple for that state fails the prior-heavy fraction bound or has too few
optional hits. Each candidate fails with probability \(o(1)\), and
these tests use independent arrays across centers.

For each odd state and each consecutive pair of levels, consider all
mappings of its neighbors into their prospective centers using at most
\(T\) IDs. Test Step~3 for each mapping, including the specified
successful-mask event when required at low states. Choose a maximal
family of failed ID sets that are pairwise disjoint, by a deterministic
local rule. Mark their union forbidden at all adjacent sites on those
levels. A site's eligible set is what remains of its prospective centers
after these marks and its singleton removals.

With conditional probability \(1-\exp(-\Omega(n))\), every site-level
ball has between \(\lambda/2\) and \(2\lambda\) prospective centers,
singleton losses are small, and every such maximal family has fewer than
\(n\) members. The first two claims follow from binomial tails.
For the third, fix presence within the count bounds. Disjoint ID sets
read disjoint arrays, which are independent given the key history.
The history's conditional failure bounds and the actual configuration
counts therefore bound the probability of \(n\) disjoint failures by
\(\exp(-\Omega(nk'_j))\) at a low star and
\(\exp(-\Omega(ns))\) at a high star, including the union over their
records. These bounds also sum over all stars and level pairs.
Since a site is incident to \(O(n)\) stars, the marks remove at most
\(O(n^2T)\) centers per site-level. The eligible sets thus have size
at least \(\lambda/3\) everywhere, with the stated exception.

Activations have not been used in defining eligibility. Apply the long
height rule and uniform eligible-active tie choices. With probability
\(1-o(1)\), neighboring even states around any odd state use at most
two consecutive levels. At either level all their choices lie in one
\((r+D)\)-ball, so the crowd bound limits their total number of IDs to
\(2n^b\le T\). Every actual mapping passes Step~3: otherwise its failed
ID set would be disjoint from the marked maximal family, contradicting
maximality. The singleton removals ensure its required hits.

\paragraph{Local validity under replacement.}
Use the same deterministic marking and singleton rules also at
hypothetical key and array values. Markings use only the displayed
likelihood tests and their records, not the history alarms or the odd
selection adjustment below. Long and short height computations use the
same pre-activation eligible sets and ties. Either returns no choice if
its maximum is \(H\) or if its resulting site-level is bad.

For either computation, call an odd state locally valid when all the
following hold at its star:
\begin{itemize}
\item every adjacent even state has a legitimate choice, and the
\(r\)-ball centered at each such state has at most \(2\lambda\)
prospective centers at every level \(0,\ldots,H\);
\item the mapping uses at most \(T\) IDs, and all required singleton
and hit tests pass;
\item the true-target gate, support conditions, and Step~3 numerical
and path tests for this mapping pass.
\end{itemize}
This is the same function of the local choices and star data for both
truncations. It does not require global path correctness or eligibility
sizes in a larger region. Thus equality of the adjacent path maxima,
with the common ties, gives equality of choices and validity. The long
calculation is valid everywhere on the global geometry event above.
Set odd kernels to zero on local invalidity, with the fixed local
fallbacks used only in auxiliary computations.

In particular, hypothetical computations do not retest global Step~1
or~2 success. The optional-prior comparisons in a high
\(\mathcal G_K\) do not read realized low-key values. Eligibility reads
optional values only in types or mappings at the consulted sites and
incident stars. Marks at a site-level do not test eligibility at other
sites. These facts ensure that the two height computations really use
the same eligible sets and that unconsulted keys remain unobserved.

\paragraph{Low rows: correcting for the selection event.}
The following rows are attached to actual odd roles, even when several
roles use one state computation. We continue to write \(b\) for the
role when taking loads, with its state understood. A high role uses the
law \(\mathcal R\) for its actual observation record. A low role
requires an adjustment because its observed arrays were selected using
the hidden column.

Fix the key history except for its target \(U_\ell\), and fix
prospective center presence. Use the original prior \(\pi_\ell\) for
the target, and simulate the raw center, activation and tie experiment
with the short height rule. Let \(\mathfrak d\) be a canonical
observation record, \(o\) its observed array data, and
\(Q_{\mathfrak d}\) the product of its Step~3 block reference laws.
When a high pool is present, include its target-independent law in this
reference. For target \(U_\ell=y\), its gated observation likelihood
is \(F_y(o)=\Theta(y)\prod_e g_e(y)\).

Let \(\mathbb P_y\) denote this raw experiment with the target replaced
by \(y\). Conditional on the recorded observations, the probability
that short selection validly presents exactly this record is a number
\(a_y(o)\in[0,1]\). More precisely,
\[
 \mathbb P_y\bigl(\text{valid proxy presentation of }\mathfrak d,
                  \ O_{\mathfrak d}\in do\bigr)
       =F_y(o)a_y(o)\,Q_{\mathfrak d}(do).
\]
The presentation event implies the required ratio gates and any successful
mask. Given these gates, the recorded low arrays have their product block
laws, independently of a recorded high pool. Integrating the remaining
selection event defines \(a_y\). In doing so, unrecorded arrays whose
laws depend on \(y\) are generated at that candidate value; they are
not fixed at their realized values. Mappings with the same record are
combined into one presentation event, and the events for different
records are disjoint. Values on zero-mass conditioning data can be fixed
arbitrarily. The resulting table \((a_y)_y\) is a function of the fixed
inputs and observed data, not a new observation of the actual target.

Set \(\varepsilon=e^{-\delta k'_j}\). At valid data define
\[
 \bar a=\frac{\sum_y\pi_\ell(y)F_ya_y}{\sum_y\pi_\ell(y)F_y}.
\]
If \(\bar a\ge\varepsilon\), use the adjusted posterior proportional
to \(\pi_\ell F a\); otherwise retain \(\mathcal P\), proportional
to \(\pi_\ell F\). For observations presented by the long rule use
this same short-rule table. In either case the row is at most
\(\varepsilon^{-1}\mathcal P\), preserving the deletion multiplier
\(e^{a_4k_c}\) and cap \(e^{D_L}/N\).

Let \(p_b^{\rm pr}\) be the proxy row, extended by zero on invalidity.
For each fixed output label \(x\), integrate the original target prior
and the raw center experiment. Cancellation of the adjusted normalizer
against the actual presentation density gives
\[
 \mathbb E[p_b^{\rm pr}(x)1_{\{\mathrm{adjusted}\}}]
 \le\pi_\ell(x)\sum_{\mathfrak d}
       \int F_x(o)a_x(o)\,dQ_{\mathfrak d}(o)
 \le\pi_\ell(x).
\]
The last step uses the disjoint presentation events. On fallback data,
its presentation density is at most \(\varepsilon\sum_y\pi_\ell(y)F_y\).
Hence
\[
 \mathbb E[p_b^{\rm pr}(x)1_{\{\mathrm{fallback}\}}]
 \le\varepsilon\pi_\ell(x)\sum_{\mathfrak d}
           \int F_x(o)\,dQ_{\mathfrak d}(o)
 \le\varepsilon(\#\mathrm{records})\pi_\ell(x).
\]
Each gated observation density integrates to at most one; their gates
need not be disjoint. The record count on the position gate is smaller
than \(\varepsilon^{-1}\), so the total mean is at most
\(2\pi_\ell(x)\).

For a fixed key history \(H\), define the normalized means
\[
 \bar p_b^{\rm pr}(y;H)=\mathbb E_{\rm raw\ centers\mid H}
                              [Np_b^{\rm pr}(y)],\qquad
 \bar p_b^{\rm long}(y;H)=\mathbb E_{\rm raw\ centers\mid H}
                              [Np_b^{\rm long}(y)].
\]
Here ``centers'' includes presence, arrays, activations and ties. Thus
presence is held fixed in the lookup calculation above and then averaged
in these means. The proxy mean uses low variables only at sign distance
\(O(\sqrt m)\) from the target, with the base and high history fixed.
Indeed the short rule reads statuses in its short tube, and the
incident-star marking rules add only bounded state distance. The one-hot
encoding bounds the corresponding sign distance. Although a consulted
center may be farther away, its array has one of these local types;
all arrays not read by the presentation event integrate out. The
validity and lookup calculations have this same locality.

At a successful key history, compare long and short calculations with
the same randomness. On correct eligibility sizes, the height lemma
bounds mismatch at any one of the \(O(n)\) neighboring states by
\(o(e^{-2m^{1/5}})\). Incorrect sizes have probability
\(e^{-\Omega(n)}\). Since \(D_L=m^{.15}\),
\[
 \bar p_b^{\rm long}(y;H)
 \le\bar p_b^{\rm pr}(y;H)
   +O(e^{D_L})\bigl[O(n)o(e^{-2m^{1/5}})+e^{-\Omega(n)}\bigr]
 =\bar p_b^{\rm pr}(y;H)+o(1).
\]
The height estimate bounds an intersection with correct sizes, not a
conditional law given size success. With either rule, the unaveraged row
and its lookup table use center randomness only within ambient radius
\(r+o(n)\) once the key history is fixed.

\paragraph{Odd column sums through the three histories.}
We now apply scattered moments three times, always averaging over actual
roles in \(B\) and setting weights to zero outside the specified class.
This distinction matters because many roles can use the same state.

First consider \(Z_b(y)=N\pi_{\ell(b)}(y)\) for low roles, under the
coarse-base law given the selected \(V_0\). The boundary roles and roles
with \(j\ge1\) contribute deterministically at most
\[
 O\left(n^{-.05+.02}+
       \sum_{h\ge1}n^{-.13h+.02(h+1)}\right)
\]
by their rarity and prior caps. For interior \(j=0\), the cap is
\(n^{.02}\). Before coarse avoidance, the mean posterior at a bin is
the bounded candidate prior, by integration of the raw conditional law.
Same-bin pairs have fraction at most \((2n^{-.04})^{d_c}\), so the
repeat cost \(n(2n^{-.04})^{d_c}n^{.02}\) tends to zero. For distinct
bins, drop the coarse constraints touching them at constant cost per
bin and integrate independent raw bin variables. The scattered-moment
bound, followed by Markov and the union over labels, shows that the
average of these weights is bounded at all labels with high probability
under the conditioned base law.

Next fix such a base and the high-key history. The weights are now
\(Z_b(y)=\bar p_b^{\rm pr}(y;H)\), viewed as functions of the low
keys. Their cap is \(e^{D_L}\). Their near relation is sign distance
\(O(\sqrt m)\), of fraction \(2^{-m+o(m)}\); hence the repeat term is
\(n2^{-m+o(m)}e^{D_L}=o(1)\). For separated rows, remove the final
low-key avoidance constraints touching their target variables. The
charge sum gives cost \(O(1)^n\). Conditional on all other low keys,
the targets are independent and no other retained row reads them, by
short locality. The one-target proxy bound gives comparison means
\(O(N\pi_{\ell(b)}(y))\); the optional pretrim adds only \(1+o(1)\)
per target. Their averages are bounded by the first calculation.
Scattered moments therefore bound all proxy-mean averages with high
probability under the low-key law. The additive comparison above gives
the same conclusion for the long means.

Finally fix one of these key histories and use weights
\(Z_b(y)=Np_b^{\rm long}(y)\), now random over raw centers. Low rows
again have cap \(e^{D_L}\). Rows separated in residual coordinates by
more than \(2r+o(n)\) have disjoint center scopes and comparison means
\(\bar p_b^{\rm long}(y;H)\). Residual-ball rarity is exponential in
\(-n\), so its product with \(ne^{D_L}\) tends to zero. High rows
contribute deterministically at most
\(2\Pr_B(j>J)e^{D_H}=o(1)\). Applying the moment and label-union
bound at this last history gives bounded normalized average loads,
while geometry success supplies valid probability rows. Since
\[
 \sum_{b\in B}p_b(y)=\frac{|B|}{N}
             \left(\frac1{|B|}\sum_{b\in B}Np_b(y)\right)
 \quad\text{and}\quad |B|/N=1/(2C_n)\longrightarrow0,
\]
all odd column sums are at most the clock threshold \(\theta_0\),
except on an event of probability \(o(1)\) intersected with the preceding
successes.

\subsection{Actual outputs and even placement}
\label{sec:actual-outputs-and-even-placement}

Fix a history and geometry with valid odd rows and the column bounds.
For each high even role \(v\), let \(c\) be its selected tuple reference.
Under the product of the odd row laws, the high neighbors' outputs
\((h_b,y_b)\) are independent. Their costs
\[
 \left[\log\frac{p_b(h_b,y_b)}
              {\widetilde{\mathcal Q}_{c,h_b}(y_b)/s}\right]_+
\]
have total mean at most \(a_4k_*n\), by the defining property of
\(\mathcal R\), and each is at most \(D_H+O(1)\) by smoothing and
the cap. Require their sum to be at most \(a_5k_*n\). Bounded-summand
concentration bounds each failure by an exponential with exponent of
order
\[
 \frac{nk_*^2}{(D_H+O(1))^2}=n^{1-.09\alpha+o(1)}.
\]
Each predicate reads at most \(n\) rows, and each row belongs to at most
\(n\) predicates. All label atoms are \(e^{o(n)}/N\), hence smaller
than any fixed inverse power of \(n\). Apply the clock lemma with
scope exponent \(B=2\). It yields an injective odd assignment, retaining
the high output indices, which satisfies all these budgets and has joint
upper bound \((1+o(1))\prod p_b\) on at most \(n^2\) queried rows.

\paragraph{Even rows: deletion of primitive block values.}
Fix an actual even role \(v\), a possible center ID there, and a tuple
reference \(c\), including its block index subset at high type. Fix the
successful key history and all primitive center data except the block
values \(z\) belonging to \(c\). Primitive data here are presence,
ties, activations and the sampled blocks themselves. In particular, do
not fix other extracted tuples that may be functions of these blocks.
The raw conditional law \(P_0'\) of \(z\) is a product of blocks of
total length \(k=k_c\), with density at most
\((A_Ke^{n^\gamma})^k\) relative to product uniform measure.

For neighboring odd outputs \(\mathbf y\), including their indices
at high rows, define the sublikelihood
\[
 F_z(\mathbf y)=1_E\prod_{b\sim v}p_b(y_b).
\]
Every choice, extracted subset and kernel is recomputed as \(z\) varies.
The local gate \(E\) requires:
\begin{itemize}
\item presence of the specified center and legitimate long-rule selection
of this \(c\) at \(v\), with its specified block subset;
\item valid neighboring rows with their count and mapping bounds, and
eligible sets of size at least \(\lambda/3\) throughout the long height
consultation domain of \(v\);
\item the prior-heavy fraction bound for \(c\), and the high log budget
at \(v\) when needed.
\end{itemize}
These conditions are local to ambient radius \(r+o(n)\), also at
hypothetical inputs. They hold for the actual selected reference on the
global success event. No success elsewhere is part of \(E\).

We construct a product reference on \(\mathbf y\) independent of
\(z\). For each same-mode neighbor \(b\), let \(\mathcal D_b(c)\) be
the finite list of possible observation configurations in its adjacent
balls under the fixed presence-count gate. Enumerate ID/type lists
without requiring that their selection is feasible for the present
array values. At a low neighbor include the possible mask; at a high
neighbor record only full pools and the possible low tuple. For each
record \(\mathfrak d\), let \(q_{b,\mathfrak d,-c}\) be the
corresponding deleted reference, namely \(\mathcal Q_c\) at low or
\(\widetilde{\mathcal Q}_{c,h}/s\) at high. Use its locally defined
fallback if its formula is undefined. Set
\[
 Q_b=\frac1{|\mathcal D_b(c)|}
          \sum_{\mathfrak d\in\mathcal D_b(c)}q_{b,\mathfrak d,-c},
 \qquad Q=\prod_{b\sim v}Q_b,
\]
using uniform \(Q_b\) at opposite-mode neighbors. If the count gate
fails, set \(E\) empty and take any fixed local reference instead.

Each deleted component uses only retained primitive data. At a low
neighbor, deletion of the full low tuple leaves any mask indicator on
a different, fixed high pool. At a high neighbor, every likelihood
factor involving a block in the fixed subset \(c\) is removed from
the one full-pool observation. The remaining target gate has no
successful-mask indicator. A pool used for several optional tuples is
still observed only once; those extracted tuples are recomputed with
\(z\), rather than held fixed. True hidden-column prefixes, sampled
before all arrays, do remain fixed. These facts prove that \(Q\) is
independent of \(z\), and that its mixture contains the actual deletion
component whenever \(E\) holds.

The likelihood costs now separate into three cases. At a low same-mode
neighbor, the adjusted row is at most \(e^{a_4k}\) times its actual
deleted component. The low record count, including a mask if needed,
fits the reserved margin because \(k'_j\le k\). At a high same-mode
neighbor use the imposed sum of positive-log costs, at most
\(a_5kn\) over these neighbors. Each full-observation-list mixture
costs only \(O(T\log n)=o(k_*)\) in the logarithm: the subset \(c\)
is fixed, and other optional subset decisions are not recorded in the
reference. Finally there are at most
\(O(mn^{.3})=o(n^{.4})\) opposite-mode neighbors, each costing
\(O(D_L+D_H)\) against uniform measure, including the high index when
present. Their total is \(o(kn)\). The reserved gaps therefore give
\[
 F_z(\mathbf y)\le e^{a_6kn}Q(\mathbf y).
\]

Let \(m_c(\mathbf y)=\int F_z(\mathbf y)\,dP_0'(z)\). For each fixed
\((v,c)\), the gated posterior lemma bounds the raw exception to
\[
 m_c(\mathbf y)>0,\qquad
 m_c(\mathbf y)\ge e^{-\delta kn}Q(\mathbf y)
\]
by \(e^{-\delta kn}\). To apply it to the actual assignment, first
retain its successful-prehistory indicator and compare the required
odd outputs to product rows using the clock bound. Only after this
comparison discard nonlocal center-success restrictions, retaining \(E\),
and integrate the raw centers at the successful key history. The bound
is thus \((1+o(1))e^{-\delta kn}\) per \((v,c)\) on the entering
success event. The union over roles, IDs and high subset indices costs
\(\exp(O(n)+O(k_*))\) and is \(o(1)\), since \(k\to\infty\).

On the complementary event use the posterior
\(F_z\,dP_0'/m_c\). Its density relative to product uniform is at most
\(e^{\tau_0kn}\): the prior contributes
\(k(\log A_K+n^\gamma)=o(kn)\), and the other exponents fit the gap
below \(\tau_0\). Let \(\bar\mu\) be its average coordinate marginal.
The support of each neighboring kernel implies that \(\bar\mu\) is
supported on common neighbors of the star.

Set \(H_1=\{x:N\bar\mu(x)>e^{\tau_1n}\}\), whose uniform mass is at
most \(e^{-\tau_1n}\). Choose
\(\upsilon_1>\tau_0/\tau_1\) with
\(\upsilon_1+\upsilon_0<1\). The joint density bound and a subset union
give
\[
 \Pr_{\rm post}\bigl(\#\{i:z_i\in H_1\}>\upsilon_1k\bigr)
 \le 2^k e^{\tau_0kn-\upsilon_1k\tau_1n}
 =e^{-\Omega(kn)}.
\]
It follows that
\[
 \bar\mu(H_1)
 =\mathbb E_{\rm post}\frac1k\sum_i1_{\{z_i\in H_1\}}
 \le\upsilon_1+e^{-\Omega(kn)}.
\]
The gate separately bounds the marginal mass on prior-heavy labels by
\(\upsilon_0\). Removing both sets leaves a mass bounded away from
zero. Normalize the restriction to obtain the even row, with normalized
cap \(O(e^{\tau_1n})\). Use only the actually selected reference, and
set the row to zero if its actual gate or its \(m_c\) test fails.

\paragraph{Even comparison means and the Hall assignment.}
We first estimate the row under product odd sampling and raw centers at
a successful key history. For a fixed potential ID and reference \(c\),
cancellation of the posterior denominator against the gated data
subdensity bounds its contribution at a prior-light label \(x\) by a
constant times
\[
 \mathbb E_{\rm centers}\!\left[
    \left(k^{-1}\sum_i1_{\{z_i=x\}}\right)
                  \sum_{\mathbf y}F_z(\mathbf y)\right].
\]
The constant accounts for the truncation normalization. The sum over
\(\mathbf y\) is at most the indicator that this center is present
and legitimately selected with the size gate. We now bound that selection
probability before averaging the tuple marginal.

Force this center present and fix its tuple values, but continue to
average all other prospective presence and activation variables. At
level zero, conditional on presence, the eligible-active tie bound is
\(3/\lambda\). At positive levels the forced-center height estimate
gives \(e^{-n^c}\) for some fixed \(c>0\). At each presence
configuration this estimate first takes the supremum over fixed
pre-activation eligible sets of the required sizes, then averages over
presence. Fixing the tuple may change these sets, but does not change
the position/activation product law, since arrays were generated
independently. Thus this is neither a bound at every fixed presence
configuration nor a posterior conditioned on selection.

The sum of presence probabilities over possible locations is
\(\lambda\) at each level. There are polynomially many levels, so
summing the selection bounds costs \(O(1)\) per fixed subset. A tuple
reference consists of whole iid blocks, and its unselected average
marginal is \(\bar P_K\), at most \(B_K/N\) on a prior-light label.
The number of possible high subsets is \(\exp(O(k_*))\). Consequently
the normalized comparison mean for the even row at \(v\) is at most
\[
 d_v=O\bigl(B_{K(v)}\exp(O(k_*1_{j(v)>J}))\bigr),
\]
whose deterministic average over actual even roles is bounded by the
cube-rarity calculation above.

For the final scattered-moment calculation set \(Z_v(x)=Np_v(x)\) for
the even rows and retain the actual success indicator. The cap is
\(O(e^{\tau_1n})\). A row's near set consists of residual-coordinate
distance at most \(2r+o(n)\). With
\(H_{\rm bin}(t)=-t\log t-(1-t)\log(1-t)\), its repeat cost is at most
\[
 n\exp\{-n(\log2-H_{\rm bin}(2\rho))+o(n)\}
                      O(e^{\tau_1n})=o(1).
\]
For up to \(n\) separated rows, at most \(n^2\) odd outputs are needed.
At a successful prehistory and fixed center data, apply the clock
comparison to these outputs. Then remove nonlocal success restrictions
but retain each local gate, and integrate the raw centers. The separated
rows have disjoint long-rule scopes, so their integrals factor and give
\[
 \mathbb E\!\left[1_{\rm success}\prod_{v\in U}Z_v(x)
                \,\middle|\,\text{successful key history}\right]
 \le (1+o(1))\prod_{v\in U}d_v
\]
for such a separated set \(U\), with harmless fixed factors included
in \(d_v\). Scattered moments and the label union give bounded normalized
average even loads. Since \(|A|/N\to0\), every even column sum is at
most one, except on an event of probability \(o(1)\) intersected with
the preceding successes.

The key-history, geometry, odd-load and predictive estimates together
have success probability tending to one under their successive stage
laws. Thus with positive probability the odd assignment is injective
and every even row is a probability law on its common neighborhood,
with all even column sums at most one. Hall's theorem gives distinct
even representatives. The two assignments form the required
\(G\)-cube.
\end{proof}
